\section{From affine coverings to a finite family}\label{b:finite-proof}

We prove Theorem~\ref{b:approximation}. On each coarse label square,
freeze the field to its affine approximation and use a much finer
square lattice. The affine covering captures the actual nonlinear
target point; its inverse estimate shows that only a thin collar of
extra centers is needed. We then estimate the cost one coarse square
at a time. Overlaps between different coarse squares are counted
with multiplicity.

\begin{proof}[Proof of Theorem~\ref{b:approximation}]
If $D=\varnothing$, take the empty collection. Otherwise fix $L,h>0$
as in the theorem. Compactness and the strict determinant bound give
\[
 \sigma=\min_{p\in D}\bigl(1+L^2\det DV(p)\bigr)>0.
\]
Choose $a>0$ so that the compact neighborhood
$D_a=\{p:\dist(p,D)\le a\}$ is contained in the domain of $V$.
Let $M$ bound $\|DV\|$ on $D_a$, and let $\omega(r)$ be the
supremum of $\|DV(x)-DV(y)\|$ over $x,y\in D_a$ with
$|x-y|\le r$. Then $\omega(r)\to0$ as $r\downarrow0$.
Set
\[
 C=\frac{1+LM}{\sigma}.
\]
By \eqref{b:affine-inverse}, this bounds the inverse of every frozen
map $I+sDV(p)$ for $p\in D$ and $|s|\le L$, including labels
where the derivative is square-zero or zero.

\paragraph{The coarse squares and their affine fields.}
Take all closed squares $S$ of side $\delta$ in a fixed square
grid that meet $D$, and choose $b_S\in S\cap D$. There are
finitely many such squares. For all sufficiently small $\delta$
they lie in $D_a$. Put
\[
 A_S=DV(b_S),\qquad V_S(p)=V(b_S)+A_S(p-b_S).
\]
The spectral hypotheses are used at $b_S\in D$ only; no such
conditions are needed on the surrounding neighborhood. Integrating
$DV$ along the segment from $b_S$ to $p\in S$ gives
\begin{equation}\label{b:finite-taylor}
 |V(p)-V_S(p)|\le\eta_\delta\delta,
 \qquad \eta_\delta=\sqrt2\,\omega(\sqrt2\delta)\longrightarrow0.
\end{equation}

\paragraph{A finite list valid for every height.}
For each $S$, apply Lemma~\ref{b:affine-cover} to $V_S$ with
fine spacing $\tau=\delta^2$. At a fine-grid center $l$ use its
closed square $T_l$ and the velocity $V_S(l)$. These are the
affine velocities; they need not equal $V(l)$. The infinite list
covers each transverse plane for every $|s|\le L$.

Fix $p\in D\cap S$ and $|s|\le L$. Apply that covering to
the actual target $p+sV(p)$, obtaining a center $l$ with
\[
 |p+sV(p)-l-sV_S(l)|\le\tau.
\]
Since $V_S(l)-V_S(p)=A_S(l-p)$,
\eqref{b:finite-taylor} gives
\[
 |(I+sA_S)(l-p)|\le\tau+L\eta_\delta\delta.
\]
The uniform inverse bound therefore yields
\begin{equation}\label{b:finite-retention}
 |l-p|\le R_\delta,
 \qquad R_\delta=C(\tau+L\eta_\delta\delta)=o(\delta).
\end{equation}
Retain all centers whose distance from the whole coarse square $S$
is at most $R_\delta$. This is a finite list independent of $p$
and $s$. The preceding argument proves that its cylinders cover
every segment with label in $D\cap S$.

\paragraph{Cost of the retained squares.}
Within the lattice belonging to $S$, the retained squares have
disjoint interiors. Their union at height zero lies within distance
$R_\delta+\tau$ of $S$. Enclosing this neighborhood in a square
enlarged along the original coordinate axes gives
\begin{equation}\label{b:finite-area}
 \sum_{l\ \mathrm{retained}}|T_l|
 \le(\delta+2R_\delta+2\tau)^2
 =\delta^2(1+o(1)).
\end{equation}
Every retained center satisfies
$|l-b_S|\le\sqrt2\delta+R_\delta=O(\delta)$, so
\[
 V_S(l)=V(b_S)+O(\delta),\qquad
 J_h(V_S(l))=J_h(V(b_S))+O(\delta).
\]
Here $h$ is fixed and all velocities lie in a bounded set. The
components in $J_h$ remain those of the original label coordinates,
regardless of the orientation of the fine squares. Combining this
estimate with \eqref{b:finite-area}, the weighted cost from $S$ is
at most
\begin{equation}\label{b:finite-square-cost}
 \delta^2J_h(V(b_S))+o(\delta^2).
\end{equation}
The errors are uniform in $S$. They depend on the fixed field,
$L,h,\sigma$, and the common modulus $\omega$, but not on a
continuous choice of orthonormal bases. In particular, changes
between hyperbolic, square-zero and zero derivatives cause no loss
in these bounds.

\paragraph{Summing over an arbitrary compact label set.}
There are $O(\delta^{-2})$ coarse squares, so the errors in
\eqref{b:finite-square-cost} sum to $o(1)$. Let $U_\delta$ be
their union. Then
\[
 D\subset U_\delta\subset D_{\sqrt2\delta}.
\]
The compact outer neighborhoods decrease to $D$ as their radius
decreases to zero. Continuity from above of Lebesgue measure gives
$|U_\delta\setminus D|\to0$, with no assumption on
$|\partial D|$. Uniform continuity of $J_h\circ V$ on $D_a$
and the disjoint interiors of the coarse squares imply
\[
 \sum_S\delta^2J_h(V(b_S))
 =\int_{U_\delta}J_h(V(p))\,dp+o(1)
 \longrightarrow\int_DJ_h(V(p))\,dp.
\]
Summing \eqref{b:finite-square-cost} thus gives the desired weighted
bound for sufficiently small $\delta$. This includes $|D|=0$,
when the total weighted cost tends to zero. The finitely many lists
together cover $E(D,V,L)$ at every height. Finally,
Proposition~\ref{b:projection} gives compact nondegenerate
parallelogram perpendicular bases under $F_h$ and their exact
physical areas $h|T_l|J_h(V_S(l))$.
\end{proof}

The physical area formula is unchanged by a Euclidean isometry.
For example, for any fixed $a\in\R$, the map
\[
 (s,\xi)\longmapsto(\xi_1,s,h(a-\xi_2))
\]
is $F_h$ followed by a coordinate permutation, a reflection and a
translation, so it gives the same perpendicular areas.

\begin{remark}[The smooth error estimate]\label{b:smooth-rate}
If $V$ has bounded second derivatives on the fixed neighborhood,
the Taylor remainder is $O(\delta^2)$ and the retained collar has
width $O(\delta^2)$. The weight varies by $O(\delta)$ on each
coarse square. The constructed weighted cost is consequently at most
\[
 \int_{U_\delta}J_h(V(p))\,dp+O(\delta)
 \le\int_DJ_h(V(p))\,dp+|U_\delta\setminus D|+O(\delta).
\]
For an arbitrary compact $D$, the excess measure tends to zero,
but the argument assigns no rate to that convergence.
\end{remark}
